% Generated draft: requires mathematical and source review.
\documentclass[11pt]{article}
\usepackage[margin=25mm]{geometry}
\usepackage{fontspec}
\setmainfont{Latin Modern Roman}
\usepackage{amsmath,amssymb,mathtools,unicode-math}
\setmathfont{Latin Modern Math}
\usepackage{enumitem,needspace,xurl,hyperref}
\setlength{\parindent}{0pt}
\setlength{\parskip}{5pt}
\setlength{\emergencystretch}{3em}
\newcommand{\N}{\mathbb{N}}
\newcommand{\Z}{\mathbb{Z}}
\newcommand{\Q}{\mathbb{Q}}
\newcommand{\R}{\mathbb{R}}
\newcommand{\C}{\mathbb{C}}
\newcommand{\F}{\mathbb{F}}
\newcommand{\A}{\mathbb{A}}
\newcommand{\PP}{\mathbb{P}}
\newcommand{\E}{\mathbb{E}}
\newcommand{\eps}{\varepsilon}
\newcommand{\dd}{\,\mathrm{d}}
\begin{document}
\section*{UnsolvedMath ID: 20003263}
\textbf{Source problem code:} catalogue entry
\subsection{Definitions and mathematical statement}
The extracted item begins ``Conjecture 1'' followed by a conditional sentence without its preceding setup. The supplied aid identifies it instead as AIM Conjecture 2. Its printed assertion is that, if $A$ in a split Cartan subgroup $D$ is not contained in a Weyl wall, every $C^1$-small perturbation of its action on $G/\Gamma$ is $C^\infty$-conjugate to an action defined by a continuous homomorphism $A\to D$. The full assumptions on $G$, $\Gamma$, the rank and connectedness of $A$, compactness and perturbation regularity are not supplied locally. They cannot be filled in by choosing a later rigidity theorem.

For the standard algebraic notation, $A$ acts by $\alpha_0(a)(g\Gamma)=ag\Gamma$. The requested conclusion for a perturbed action $\beta$ is existence of a smooth diffeomorphism $h$ and continuous homomorphism $\eta:A\to D$ with
\[
 h\circ\beta(a)\circ h^{-1}=L_{\eta(a)}
 \qquad\text{for every }a\in A.
\]
The phrase every sufficiently small perturbation would require a neighborhood in a specified action-space $C^1$ topology. Smoothness of the perturbed action is needed for the usual smooth-conjugacy reading; neither that regularity nor closeness requirements on $h$ or $\eta$ are stated. The normal form is specifically in $D$, not an unspecified larger centralizer.

The aid's hypothesis-separation result fixes connected $A$ with Lie algebra $\mathfrak a\subseteq\mathfrak d=\operatorname{Lie}(D)$. For each root $\alpha$ relative to the split Cartan, a wall is $\ker\alpha$. In this setting
\[
 A\text{ lies in no Weyl wall}
 \quad\Longleftrightarrow\quad
 \alpha|_{\mathfrak a}\ne0\text{ for every root }\alpha.
\]
The complement in $\mathfrak a$ of the finitely many restricted root kernels is open dense, providing regular elements, with the corresponding hyperbolic root directions and neutral centralizer directions. It is not a claim of Anosov behavior with no center or of a full local-rigidity theorem.

The explicit root-detecting plane in $\mathfrak{sl}_4(\R)$ is
\[
 \mathfrak a=\{\operatorname{diag}(s,t,-s,-t):s,t\in\R\}.
\]
Every root is nonzero on this plane, but the distinct roots $e_1-e_4$ and $e_2-e_3$ both restrict to $s+t$. Thus wall avoidance does not imply the stronger root-projective-genericity condition requiring appropriate distinct nonproportional ambient roots to remain nonproportional. The exact older theorem's formulation still needs its reference.

Separately, for two simple factors with regular Cartan elements $H_1,H_2$, the two-dimensional subgroup generated independently by $(H_1,0)$ and $(0,H_2)$ detects every root of the product while having rank-one projection to each factor. Root detection therefore does not imply higher rank in every simple-factor projection. This algebraic example is not automatically an example satisfying any omitted irreducible-lattice hypotheses of a rigidity theorem.

The summary says later rigidity theorems cover broad partial cases but do not on the verified hypotheses and normal forms settle every quantifier of the printed conjecture. No full citations, counterexample to the complete original conjecture, or general resolution are supplied. In particular, the missing rank assumptions and the numbering discrepancy must remain visible.
\subsection{Short English statement}
The aid separates wall avoidance from stronger hypotheses used in local-rigidity theorems. Its two examples detect every root while violating projective root genericity or factorwise higher rank. Those separations do not decide the full conjecture, whose preceding setup and action-topology conventions are missing.
\subsection{Sources}
\begin{itemize}
\item Local UnsolvedMath registry, permanent ID 20003263, source code catalogue entry; \texttt{aim-other-notes.json}, notes[69], source tag conjecture.
\item AIM, \emph{Emerging applications of measure rigidity}: \url{https://aimath.org/WWN/measrigid/measrigid.pdf}. The extraction says Conjecture 1; the supplied research summary identifies Conjecture 2. No full preceding setup or proposer is supplied.
\item Older projective-genericity and later local-rigidity theorems are referenced only generically in the summary, without authors, titles or precise statements.
\item Machine-generated research attempt 2: partial result, partially solved, rigor labelled ``proved,'' low novelty confidence. Not peer reviewed; independently verify before citation. Registry updated 14 August 2026.
\end{itemize}
\end{document}
